\documentclass{article}
\usepackage{amsmath, amssymb}

\title{Rigid-Body Equations of Motion from a Discrete-Mass Model}
\author{}
\date{}

\begin{document}
\maketitle

\section{Kinematic representation}

Consider a rigid body represented by discrete masses \(m_i\), with inertial-frame positions

\[
x_i(t) = b(t) + A(t)y_i,
\]

where

\[
b(t) \in \mathbb{R}^3
\]

is the location of the center of mass,

\[
y_i \in \mathbb{R}^3
\]

are fixed coordinates attached to the body, and

\[
A(t) \in SO(3)
\]

is the rotation matrix mapping body-frame vectors into the inertial frame.

Since \(A\) is a rotation matrix,

\[
A^T A = I.
\]

Choose the body-frame origin to coincide with the center of mass. Then

\[
\sum_i m_i y_i = 0.
\]

Define the total mass

\[
M = \sum_i m_i.
\]

\section{Discrete equations of motion}

Let \(f_i\) denote the external force acting on mass \(i\), and let \(g_{ij}\) denote the internal force exerted by mass \(j\) on mass \(i\).

The total force on mass \(i\) is

\[
F_i = f_i+\sum_{j\ne i}g_{ij}.
\]

Newton's second law gives

\[
m_i\ddot{x}_i = f_i+\sum_{j\ne i}g_{ij}.
\]

From

\[
x_i=b+Ay_i,
\]

and because \(y_i\) is constant,

\[
\dot{x}_i = \dot{b}+\dot{A}y_i,
\]

and

\[
\ddot{x}_i = \ddot{b}+\ddot{A}y_i.
\]

Thus

\[
m_i\left(\ddot{b}+\ddot{A}y_i\right) = f_i+\sum_{j\ne i}g_{ij}.
\]

\section{Translational equation}

Sum the equations over all masses:

\[
\sum_i m_i\ddot{b} + \sum_i m_i\ddot{A}y_i = \sum_i f_i + \sum_i\sum_{j\ne i}g_{ij}.
\]

Since \(\ddot b\) and \(\ddot A\) are independent of \(i\),

\[
M\ddot b + \ddot A\sum_i m_i y_i = \sum_i f_i + \sum_i\sum_{j\ne i}g_{ij}.
\]

Using

\[
\sum_i m_i y_i=0,
\]

and Newton's third law,

\[
g_{ij}=-g_{ji},
\]

the internal forces cancel pairwise:

\[
\sum_i\sum_{j\ne i}g_{ij}=0.
\]

Therefore,

\[
M\ddot b=\sum_i f_i.
\]

Hence

\[
\ddot b = \frac{1}{M}\sum_i f_i.
\]

\section{Rotational equation in matrix form}

Return to the individual equations,

\[
m_i\left(\ddot b+\ddot A y_i\right) = f_i+\sum_{j\ne i}g_{ij}.
\]

Multiply on the right by \(y_i^T\) and sum over \(i\):

\[
\sum_i m_i\ddot b\,y_i^T + \sum_i m_i\ddot A y_i y_i^T = \sum_i f_i y_i^T + \sum_i\sum_{j\ne i}g_{ij}y_i^T.
\]

The first term vanishes because

\[
\sum_i m_i y_i=0.
\]

Define the body-frame second-moment matrix

\[
C = \sum_i m_i y_i y_i^T.
\]

Then

\[
\ddot A C = \sum_i f_i y_i^T + \sum_i\sum_{j\ne i}g_{ij}y_i^T.
\]

Assuming \(C\) is nonsingular,

\[
\ddot A = \left[ \sum_i f_i y_i^T + \sum_i\sum_{j\ne i}g_{ij}y_i^T \right]C^{-1}.
\]

This is the direct equation obtained from the discrete particle equations. However, it does not yet explicitly enforce the fact that \(A\) must remain a rotation matrix. That constraint is obtained from \(A^T A=I\).

\section{Consequences of \((A^T A)'=0\)}

Because

\[
A^T A=I,
\]

differentiate with respect to time:

\[
\left(A^T A\right)' = \dot A^T A+A^T\dot A = 0.
\]

Therefore,

\[
\dot A^T A+A^T\dot A=0.
\]

Define

\[
\Omega=A^T\dot A.
\]

Then

\[
\Omega^T = \dot A^T A = -A^T\dot A = -\Omega.
\]

Thus

\[
\Omega^T=-\Omega.
\]

Therefore \(\Omega\) is skew-symmetric and contains only three independent components.

Since

\[
A^TA=I,
\]

we can multiply the definition of \(\Omega\) by \(A\) on the left:

\[
A\Omega = AA^T\dot A = \dot A.
\]

Hence

\[
\dot A=A\Omega.
\]

The skew-symmetric matrix \(\Omega\) is the body-frame angular-velocity matrix. It may equivalently be written as

\[
\Omega=\widehat{\omega},
\]

where

\[
\widehat{\omega} = \begin{bmatrix} 0&-\omega_3&\omega_2\\ \omega_3&0&-\omega_1\\ -\omega_2&\omega_1&0 \end{bmatrix}.
\]

\section{Second derivative of the orthogonality constraint}

Differentiate

\[
\dot A^T A+A^T\dot A=0
\]

again:

\[
\ddot A^T A + \dot A^T\dot A + \dot A^T\dot A + A^T\ddot A = 0.
\]

Therefore,

\[
\ddot A^T A + 2\dot A^T\dot A + A^T\ddot A = 0.
\]

Equivalently,

\[
A^T\ddot A + \ddot A^T A = -2\dot A^T\dot A.
\]

Thus the symmetric part of \(A^T\ddot A\) is constrained by \(\dot A\):

\[
\operatorname{sym}(A^T\ddot A) = -\dot A^T\dot A.
\]

Using

\[
\dot A=A\Omega,
\]

we obtain

\[
\dot A^T\dot A = \Omega^T A^T A\Omega = \Omega^T\Omega.
\]

Since

\[
\Omega^T=-\Omega,
\]

we have

\[
\Omega^T\Omega = -\Omega^2.
\]

Therefore,

\[
\operatorname{sym}(A^T\ddot A) = \Omega^2.
\]

\section{Decomposition of \(A^T\ddot A\)}

Differentiate

\[
\dot A=A\Omega:
\]

\[
\ddot A = \dot A\Omega + A\dot\Omega.
\]

Substitute

\[
\dot A=A\Omega:
\]

\[
\ddot A = A\Omega^2+A\dot\Omega.
\]

Hence

\[
\ddot A = A\left(\dot\Omega+\Omega^2\right).
\]

Multiplying on the left by \(A^T\),

\[
A^T\ddot A = \dot\Omega+\Omega^2.
\]

Now,

\[
\dot\Omega^T = -\dot\Omega,
\]

so \(\dot\Omega\) is skew-symmetric, while

\[
(\Omega^2)^T = (\Omega^T)^2 = (-\Omega)^2 = \Omega^2,
\]

so \(\Omega^2\) is symmetric.

Consequently,

\[
\operatorname{skew}(A^T\ddot A)=\dot\Omega,
\]

and

\[
\operatorname{sym}(A^T\ddot A)=\Omega^2.
\]

This is the key simplification obtained from the constraint \(A^T A=I\).

\section{Apply the constraint to the force equation}

Define

\[
Q = \sum_i f_i y_i^T + \sum_i\sum_{j\ne i}g_{ij}y_i^T.
\]

Then the discrete equations give

\[
\ddot A C=Q,
\]

or

\[
\ddot A=QC^{-1}.
\]

Therefore,

\[
A^T\ddot A = A^TQC^{-1}.
\]

Using

\[
A^T\ddot A = \dot\Omega+\Omega^2,
\]

we obtain

\[
\dot\Omega+\Omega^2 = A^TQC^{-1}.
\]

Taking the skew-symmetric part gives the rotational evolution equation:

\[
\dot\Omega = \operatorname{skew} \left(A^TQC^{-1}\right).
\]

Thus,

\[
\dot\Omega = \operatorname{skew} \left[ A^T \left( \sum_i f_i y_i^T + \sum_i\sum_{j\ne i}g_{ij}y_i^T \right) C^{-1} \right].
\]

This equation contains only the three independent rotational degrees of freedom.

\section{Simplification of the internal-force term}

Assume the internal forces obey Newton's third law:

\[
g_{ij}=-g_{ji}.
\]

The internal contribution is

\[
\sum_i\sum_{j\ne i}g_{ij}y_i^T.
\]

Group the terms belonging to each pair \(i,j\). For one pair,

\[
g_{ij}y_i^T+g_{ji}y_j^T = g_{ij}(y_i-y_j)^T.
\]

Therefore,

\[
\sum_i\sum_{j\ne i}g_{ij}y_i^T = \sum_{i<j}g_{ij}(y_i-y_j)^T.
\]

Hence

\[
Q = \sum_i f_i y_i^T + \sum_{i<j}g_{ij}(y_i-y_j)^T.
\]

\section{Central internal forces}

For ordinary rigid-body constraint forces, assume each internal force is directed along the line connecting the two masses:

\[
g_{ij} = \lambda_{ij}(x_j-x_i).
\]

Since

\[
x_j-x_i = A(y_j-y_i),
\]

we have

\[
g_{ij} = \lambda_{ij}A(y_j-y_i).
\]

Let

\[
q_{ij}=y_j-y_i.
\]

Then

\[
g_{ij} = \lambda_{ij}Aq_{ij}.
\]

The pair contribution to \(Q\) is

\[
g_{ij}(y_i-y_j)^T = -\lambda_{ij}Aq_{ij}q_{ij}^T.
\]

Therefore,

\[
A^Tg_{ij}(y_i-y_j)^T = -\lambda_{ij}q_{ij}q_{ij}^T.
\]

But

\[
q_{ij}q_{ij}^T
\]

is symmetric. Hence

\[
\operatorname{skew} \left[ A^Tg_{ij}(y_i-y_j)^T \right] = 0.
\]

Therefore all central internal pair forces disappear from the skew part of the rotational equation.

Thus,

\[
\dot\Omega = \operatorname{skew} \left[ A^T \left( \sum_i f_i y_i^T \right) C^{-1} \right].
\]

This is the rotational equation when the internal forces are central constraint forces satisfying Newton's third law.

\section{What happens to the internal forces?}

The internal forces have not simply been discarded.

The complete equation is

\[
\ddot A = \left[ \sum_i f_i y_i^T + \sum_{i<j}g_{ij}(y_i-y_j)^T \right]C^{-1}.
\]

The internal forces contribute to the symmetric part of

\[
A^T\ddot A.
\]

But the symmetric part is already determined kinematically:

\[
\operatorname{sym}(A^T\ddot A)=\Omega^2.
\]

Thus the internal constraint forces enforce

\[
\operatorname{sym}(A^T\ddot A)=\Omega^2,
\]

while the external forces determine the skew part,

\[
\operatorname{skew}(A^T\ddot A)=\dot\Omega.
\]

This gives the decomposition

\[
A^T\ddot A = \underbrace{\dot\Omega}_{\text{rotational dynamics}} + \underbrace{\Omega^2}_{\text{rigid-body kinematics}}.
\]

\section{Final equations for numerical integration}

The system can therefore be integrated using the state variables

\[
b,\qquad v=\dot b,\qquad A,\qquad \Omega.
\]

The equations are

\[
\dot b=v
\]

\[
\dot v = \frac{1}{M}\sum_i f_i
\]

\[
\dot A=A\Omega
\]

and, for central internal forces,

\[
\dot\Omega = \operatorname{skew} \left[ A^T \left( \sum_i f_i y_i^T \right) C^{-1} \right].
\]

where

\[
C=\sum_i m_i y_i y_i^T.
\]

The particle positions can then be recovered from

\[
x_i=b+Ay_i.
\]

Their accelerations are

\[
\ddot x_i = \ddot b+\ddot A y_i,
\]

with

\[
\ddot A = A(\dot\Omega+\Omega^2).
\]

\section{Relation to the conventional inertia tensor}

The matrix \(C\) is not itself the conventional inertia tensor. The body-frame inertia tensor is

\[
I = \sum_i m_i \left[ (y_i^Ty_i)I_3-y_i y_i^T \right].
\]

Using

\[
\operatorname{tr}(C) = \sum_i m_i y_i^Ty_i,
\]

we can write

\[
I = \operatorname{tr}(C)I_3-C.
\]

Thus \(C\) and the usual inertia tensor contain the same information.

If

\[
\Omega=\widehat{\omega},
\]

then the matrix equation above is equivalent to Euler's rigid-body equation

\[
I\dot\omega+\omega\times(I\omega)=\tau,
\]

where the body-frame torque is

\[
\tau = A^T \sum_i(Ay_i)\times f_i.
\]

The matrix formulation is therefore not a different physical model;
it is the same rigid-body dynamics expressed directly in terms of the rotation matrix and the discrete-mass representation.


\section{Final note}

A useful implementation detail is that I would integrate \(A\) and \(\Omega\), rather than \(\ddot A\).
The constraint \(A^TA=I\) then reduces the rotational state to three angular-velocity variables,
and you can enforce \(A\in SO(3)\) numerically using a rotation-preserving update (or quaternions)
rather than allowing nine independently integrated entries of \(A\) to drift.
\end{document}
